\documentclass[a4,11pt]{aleph-notas}
% Se puede ver la documentación aquí:
% https://github.com/alephsub0/LaTeX_aleph-notas

% -- Paquetes adicionales
\usepackage{enumitem}
\usepackage{url}
\usepackage{array}
\usepackage{booktabs}
\usepackage{longtable}
\usepackage{aleph-comandos}
\hypersetup{
    urlcolor=blue,
    linkcolor=blue,
}


% -- Datos
\institucion{Facultad de Ciencias Exactas, Naturales y Ambientales}
\carrera{Catálogo STEM}
\asignatura{Matemática III}
\tema{Plan tema 17: Teoremas integrales}
\autor{Andrés Merino}
\fecha{Periodo 2026-1}

\logouno[0.16\textwidth]{Logos/logoPUCE_04_ac}
\definecolor{colortext}{HTML}{1957d1}
\definecolor{colordef}{HTML}{1957d1}
\fuente{montserrat}


\begin{document}

\encabezado

%%%%%%%%%%%%%%%%%%%%%%%%%%%%%%%%%%%%%%%%
\section*{Resultado de Aprendizaje}
%%%%%%%%%%%%%%%%%%%%%%%%%%%%%%%%%%%%%%%%

%%%%%%%%%%%%%%%%%%%%%%%%%%%%%%%%%%%%%%%%
\subsection*{RdA de la asignatura:}
\begin{itemize}[leftmargin=*]
    \item \textbf{RdA 1:} Identificar conceptos, teoremas y propiedades de funciones vectoriales, derivadas parciales e integral vectorial.
    \item \textbf{RdA 2:} Calcular derivadas parciales e integrales con asistentes computacionales.
    \item \textbf{RdA 3:} Aplicar Cálculo Vectorial para resolver problemas reales.
\end{itemize}

%%%%%%%%%%%%%%%%%%%%%%%%%%%%%%%%%%%%%%%%
\subsection*{RdA del tema:}
\begin{itemize}[leftmargin=*]
    \item Aplicar el Teorema de Green para relacionar una integral de línea sobre una curva cerrada con una integral doble.
    \item Aplicar el Teorema de Gauss en el plano para calcular el flujo a través de una curva cerrada.
    \item Aplicar el Teorema de Stokes para relacionar la integral de superficie del rotacional con una integral de línea.
    \item Aplicar el Teorema de Gauss (de la divergencia) para relacionar el flujo a través de una superficie cerrada con una integral triple.
\end{itemize}

%%%%%%%%%%%%%%%%%%%%%%%%%%%%%%%%%%%%%%%%
\section*{Introducción}
%%%%%%%%%%%%%%%%%%%%%%%%%%%%%%%%%%%%%%%%

\paragraph{Pregunta inicial:}
Para saber cuánto aire entra y sale de una habitación, ¿será necesario medir el flujo en cada punto de las paredes, o bastará con conocer lo que ocurre en el interior? En general, ¿podemos deducir lo que sucede en el borde de una región a partir de lo que pasa dentro de ella?

%%%%%%%%%%%%%%%%%%%%%%%%%%%%%%%%%%%%%%%%
\section*{Desarrollo}
%%%%%%%%%%%%%%%%%%%%%%%%%%%%%%%%%%%%%%%%

%%%%%%%%%%%%%%%%%%%%%%%%%%%%%%%%%%%%%%%%
\subsection*{Actividad 1: Teoremas de Green y de Gauss en el plano}
Se presentan los teoremas integrales en el plano, que relacionan una integral de línea sobre una curva cerrada con una integral doble sobre la región que encierra.

\paragraph{¿Cómo lo haremos?}
\begin{itemize}[leftmargin=*]
    \item \textbf{Motivación:} retomar la pregunta inicial y mostrar que, en el plano, la circulación y el flujo a lo largo del borde de una región pueden calcularse integrando una cantidad en su interior.
    \item \textbf{Clase magistral:} se enuncia el Teorema de Green en sus formas $\iint_\Omega \rot(f)\,dA=\oint_{\partial\Omega} f\cdot ds$ y $\iint_\Omega(\partial N/\partial x-\partial M/\partial y)\,dxdy=\oint_{\partial\Omega} M\,dx+N\,dy$; se enuncia el Teorema de Gauss en el plano para el flujo a través del borde. Se usa el resumen \href{https://andres-merino.github.io/MatematicaIII-05-N0051/01Resumen/Resumen17.pdf}{Resumen17.pdf}.
    \item \textbf{Implementación computacional:} verificar el Teorema de Green calculando ambos lados de la igualdad con un asistente computacional.
    \item \textbf{Resolución de ejercicios:} aplicar los teoremas de Green y de Gauss en el plano usando \href{https://andres-merino.github.io/MatematicaIII-05-N0051/04Ejercicios/Ejercicios17.pdf}{Ejercicios17.pdf}.
\end{itemize}

\paragraph{Verificación de aprendizaje:}
\begin{itemize}[leftmargin=*]
    \item Use el Teorema de Green para calcular $\oint_{\partial\Omega} -y\,dx+x\,dy$ donde $\partial\Omega$ es la circunferencia unitaria.
    \item ¿Qué hipótesis debe cumplir la región $\Omega$ y su borde para poder aplicar el Teorema de Green?
    \item ¿Qué relación tiene el Teorema de Green con el cálculo de áreas planas?
\end{itemize}

%%%%%%%%%%%%%%%%%%%%%%%%%%%%%%%%%%%%%%%%
\subsection*{Actividad 2: Teoremas de Stokes y de Gauss}
Se generalizan los teoremas integrales al espacio: el Teorema de Stokes relaciona el rotacional con una integral de línea y el Teorema de Gauss relaciona la divergencia con el flujo a través de una superficie cerrada.

\paragraph{¿Cómo lo haremos?}
\begin{itemize}[leftmargin=*]
    \item \textbf{Motivación:} presentar el Teorema de Stokes como la versión espacial del Teorema de Green y el Teorema de Gauss como el balance entre lo que ocurre dentro de un sólido y el flujo a través de su frontera.
    \item \textbf{Clase magistral:} se introducen las superficies orientables y la orientación consistente del borde; se enuncia el Teorema de Stokes $\iint_\Sigma \rot(f)\cdot dS=\oint_{\partial\Sigma} f\cdot ds$ y el Teorema de Gauss $\iiint_\Omega \dive(F)\,dV=\varoiint_{\partial\Omega} F\cdot dS$. Se usa el resumen \href{https://andres-merino.github.io/MatematicaIII-05-N0051/01Resumen/Resumen17.pdf}{Resumen17.pdf}.
    \item \textbf{Resolución de ejercicios:} aplicar los teoremas de Stokes y de Gauss usando \href{https://andres-merino.github.io/MatematicaIII-05-N0051/04Ejercicios/Ejercicios17.pdf}{Ejercicios17.pdf}.
\end{itemize}

\paragraph{Verificación de aprendizaje:}
\begin{itemize}[leftmargin=*]
    \item Use el Teorema de Gauss para calcular el flujo de $F(x,y,z)=(x,y,z)$ a través de la esfera unitaria.
    \item ¿En qué se parecen y en qué se diferencian los teoremas de Green y de Stokes?
    \item Identifique el patrón común a los teoremas de Green, Stokes y Gauss: ¿qué relacionan en cada caso el interior de una región y su frontera?
\end{itemize}

%%%%%%%%%%%%%%%%%%%%%%%%%%%%%%%%%%%%%%%%
\section*{Cierre}
%%%%%%%%%%%%%%%%%%%%%%%%%%%%%%%%%%%%%%%%

\paragraph{Tarea:}
Resolver del libro \href{https://catalogobiblioteca.puce.edu.ec/cgi-bin/koha/opac-detail.pl?biblionumber=83405&query_desc=su%3A%22An%C3%A1lisis%20vectorial%22}{Cálculo Vectorial de Colley}: sección 6.2, ejercicios 1--15 (impares).

\paragraph{Pregunta de investigación:}
\begin{enumerate}[leftmargin=*]
    \item ¿Cómo se unifican los teoremas de Green, Stokes y Gauss en el Teorema General de Stokes mediante formas diferenciales?
    \item ¿Qué papel juegan estos teoremas en las ecuaciones de Maxwell del electromagnetismo? Investigar un ejemplo concreto.
\end{enumerate}

\paragraph{Para el próximo tema:}
Repasar de manera integral los contenidos del curso, articulando funciones vectoriales, derivadas parciales e integrales múltiples y vectoriales como un todo coherente, en preparación para la evaluación final.


\end{document}
